\documentclass[preview]{standalone}

\usepackage{amsmath}
\usepackage{amssymb}
\usepackage{stellar}
\usepackage{definitions}
\usepackage{bettelini}

\begin{document}

\id{probability-discrete-random-vectors}
\genpage

\section{Discrete random vectors}

\begin{snippetdefinition}{discrete-random-vector-definition}{Discrete random vector}
    Let \((\Omega,\powerset(\Omega),\mathbb P)\) be a
    \countableprobabilityspace, and let
    \(n\in\naturalnumbers^\exceptzero\). For every \(i\in\{1,\ldots,n\}\),
    let \(d_i\in\naturalnumbers^\exceptzero\), and let
    \(E_i\subseteq\realnumbers^{d_i}\) be a finite or countably infinite
    \set.
    A \emph{discrete random vector} is an \(n\)-tuple
    \[
        X=(X_1,\ldots,X_n)
    \]
    such that each \(X_i\colon\Omega\fromto E_i\) is a
    \discreterandomvariable. Equivalently, \(X\) is the
    \discreterandomvariable, identifying \(\prod_{i=1}^n E_i\) with a
    finite or countably infinite subset of
    \(\realnumbers^{d_1+\cdots+d_n}\),
    \[
        X\colon\Omega\fromto\prod_{i=1}^n E_i,
        \qquad
        X(\omega)=\bigl(X_1(\omega),\ldots,X_n(\omega)\bigr).
    \]
\end{snippetdefinition}

\begin{snippetdefinition}{joint-law-density-definition}{Joint law and joint density}
    Let \((\Omega,\powerset(\Omega),\mathbb P)\) be a
    \countableprobabilityspace, and let
    \(X=(X_1,\ldots,X_n)\) be a \discreterandomvector with
    \(X_i\colon\Omega\fromto E_i\). The \emph{joint law} of \(X\) is the
    \lawofrandomvariable of
    \[
        X\colon\Omega\fromto\prod_{i=1}^n E_i.
    \]
    It is denoted by \(\mu_{X_1,\ldots,X_n}\).
    The \emph{joint density} of \(X\) is the \randomvariabledensity
    \[
        p_{X_1,\ldots,X_n}\colon\prod_{i=1}^n E_i\fromto[0,1]
    \]
    defined by
    \[
        p_{X_1,\ldots,X_n}(x_1,\ldots,x_n)
        =
        \mathbb P(X_1=x_1,\ldots,X_n=x_n).
    \]
\end{snippetdefinition}

\begin{snippetproposition}{joint-law-on-rectangles}{Joint law on rectangles}
    Let \((\Omega,\powerset(\Omega),\mathbb P)\) be a
    \countableprobabilityspace, and let
    \(X=(X_1,\ldots,X_n)\) be a \discreterandomvector with
    \(X_i\colon\Omega\fromto E_i\). For every
    \(A_i\subseteq E_i\),
    \[
        \mu_{X_1,\ldots,X_n}\left(\prod_{i=1}^n A_i\right)
        =
        \mathbb P(X_1\in A_1,\ldots,X_n\in A_n).
    \]
\end{snippetproposition}

\begin{snippetproof}{joint-law-on-rectangles-proof}{joint-law-on-rectangles}{Joint law on rectangles}
    By definition of the \jointlaw,
    \[
        \mu_{X_1,\ldots,X_n}\left(\prod_{i=1}^n A_i\right)
        =
        \mathbb P\left(
            X\in\prod_{i=1}^n A_i
        \right).
    \]
    Since \(X(\omega)=(X_1(\omega),\ldots,X_n(\omega))\), the \event
    \(X\in\prod_{i=1}^n A_i\) is exactly the \event
    \(X_1\in A_1,\ldots,X_n\in A_n\).
\end{snippetproof}

\section{Marginals}

\begin{snippetdefinition}{marginal-law-density-definition}{Marginal law and marginal density}
    Let \((\Omega,\powerset(\Omega),\mathbb P)\) be a
    \countableprobabilityspace, and let
    \(X=(X_1,\ldots,X_n)\) be a \discreterandomvector with
    \(X_i\colon\Omega\fromto E_i\). Let
    \[
        I=(i_1,\ldots,i_k)
    \]
    with \(1\leq i_1<\cdots<i_k\leq n\).
    The \emph{marginal law} of \(X\) along \(I\) is the
    \lawofrandomvariable of the subvector
    \[
        X_I=(X_{i_1},\ldots,X_{i_k})
        \colon\Omega\fromto\prod_{j=1}^k E_{i_j}.
    \]
    The \emph{marginal density} along \(I\) is the
    \randomvariabledensity of \(X_I\), denoted by \(p_I\).
\end{snippetdefinition}

\begin{snippetproposition}{marginal-density-from-joint-density}{Marginal density from the joint density}
    Let \((\Omega,\powerset(\Omega),\mathbb P)\) be a
    \countableprobabilityspace, and let
    \(X=(X_1,\ldots,X_n)\) be a \discreterandomvector with
    \(X_i\colon\Omega\fromto E_i\). Let
    \(I=(i_1,\ldots,i_k)\) with \(1\leq i_1<\cdots<i_k\leq n\), and let
    \(p_I\) be the \marginaldensity along \(I\). Then
    \[
        p_I(x_{i_1},\ldots,x_{i_k})
        =
        \sum_{\substack{x_j\in E_j\\j\notin\{i_1,\ldots,i_k\}}}
        p_{X_1,\ldots,X_n}(x_1,\ldots,x_n).
    \]
\end{snippetproposition}

\begin{snippetproof}{marginal-density-from-joint-density-proof}{marginal-density-from-joint-density}{Marginal density from the joint density}
    The \event
    \[
        \{X_{i_1}=x_{i_1},\ldots,X_{i_k}=x_{i_k}\}
    \]
    is the disjoint union, over all choices \(x_j\in E_j\) with
    \(j\notin\{i_1,\ldots,i_k\}\), of the \event[events]
    \[
        \{X_1=x_1,\ldots,X_n=x_n\}.
    \]
    By countable additivity,
    \[
        \mathbb P(X_{i_1}=x_{i_1},\ldots,X_{i_k}=x_{i_k})
        =
        \sum_{\substack{x_j\in E_j\\j\notin\{i_1,\ldots,i_k\}}}
        \mathbb P(X_1=x_1,\ldots,X_n=x_n),
    \]
    which is the desired formula.
\end{snippetproof}

\begin{snippetcorollary}{two-dimensional-marginal-density-formulas}{Two-dimensional marginal densities}
    Let \((\Omega,\powerset(\Omega),\mathbb P)\) be a
    \countableprobabilityspace, let
    \(X\colon\Omega\fromto E_1\) and \(Y\colon\Omega\fromto E_2\) be
    \discreterandomvariable[discrete random variables], and let \(p_{X,Y}\)
    be their \jointdensity. Then
    \[
        p_X(x)=\sum_{y\in E_2}p_{X,Y}(x,y),
        \qquad
        p_Y(y)=\sum_{x\in E_1}p_{X,Y}(x,y).
    \]
    Moreover, for every \(A_1\subseteq E_1\) and \(A_2\subseteq E_2\),
    \[
        \mu_{X,Y}(A_1\cartesianprod A_2)
        =
        \sum_{x\in A_1}\sum_{y\in A_2}p_{X,Y}(x,y).
    \]
\end{snippetcorollary}

\begin{snippetproof}{two-dimensional-marginal-density-formulas-proof}{two-dimensional-marginal-density-formulas}{Two-dimensional marginal densities}
    The first two formulas are the special case \(n=2\) of the
    marginalization formula. The formula for \(A_1\cartesianprod A_2\)
    follows from the \jointdensity:
    \[
        \mu_{X,Y}(A_1\cartesianprod A_2)
        =
        \sum_{(x,y)\in A_1\cartesianprod A_2}p_{X,Y}(x,y)
        =
        \sum_{x\in A_1}\sum_{y\in A_2}p_{X,Y}(x,y).
    \]
\end{snippetproof}

\begin{snippetproposition}{zero-marginal-implies-zero-joint-slice}{Zero marginal slices}
    Let \((\Omega,\powerset(\Omega),\mathbb P)\) be a
    \countableprobabilityspace, let
    \(X\colon\Omega\fromto E_1\) and \(Y\colon\Omega\fromto E_2\) be
    \discreterandomvariable[discrete random variables], and let \(p_{X,Y}\)
    be their \jointdensity. If \(p_X(x)=0\), then
    \[
        p_{X,Y}(x,y)=0
    \]
    for every \(y\in E_2\). If \(p_Y(y)=0\), then
    \[
        p_{X,Y}(x,y)=0
    \]
    for every \(x\in E_1\).
\end{snippetproposition}

\begin{snippetproof}{zero-marginal-implies-zero-joint-slice-proof}{zero-marginal-implies-zero-joint-slice}{Zero marginal slices}
    If \(p_X(x)=0\), then
    \[
        0=p_X(x)=\sum_{y\in E_2}p_{X,Y}(x,y).
    \]
    Every term in the sum is nonnegative, so each term must be \(0\).
    The proof for \(p_Y(y)=0\) is analogous.
\end{snippetproof}

\section{Joint laws are not determined by marginals}

\begin{snippetexample}{same-marginals-different-joint-laws-urn-example}{Same marginals but different joint laws}
    Let \(n\in\naturalnumbers\) with \(n\geq2\), and let
    \(E=\{1,\ldots,n\}\). Consider two experiments in which two balls are
    drawn from an urn containing \(n\) distinguishable balls numbered by \(E\).

    In the experiment with replacement, let
    \[
        \Omega_{\mathrm{wr}}=E^2
    \]
    with the uniform probability, and let
    \(X_1(i,j)=i\), \(X_2(i,j)=j\). Then
    \[
        p_{X_1,X_2}^{\mathrm{wr}}(i,j)=\frac{1}{n^2}
    \]
    for every \((i,j)\in E^2\).

    In the experiment without replacement, let
    \[
        \Omega_{\mathrm{wor}}=\{(i,j)\in E^2\suchthat i\neq j\}
    \]
    with the uniform probability, and let
    \(Y_1(i,j)=i\), \(Y_2(i,j)=j\). Then
    \[
        p_{Y_1,Y_2}^{\mathrm{wor}}(i,j)
        =
        \begin{cases}
            \frac{1}{n(n-1)} & i\neq j,\\
            0 & i=j.
        \end{cases}
    \]
    In both experiments, the two one-dimensional \marginallaw[marginal laws]
    are uniform on \(E\). However, the \jointlaw[joint laws] are different.
    Therefore the \marginallaw[marginal laws] do not determine the
    \jointlaw.
\end{snippetexample}

\begin{snippetexample}{random-vector-support-not-product-example}{A joint support need not be a product}
    In the sampling-without-replacement model of the previous example, the
    \randomvariablesupport of the \discreterandomvector \((Y_1,Y_2)\) is
    \[
        \{(i,j)\in E^2\suchthat i\neq j\}.
    \]
    The marginal \randomvariablesupport[supports] of \(Y_1\) and \(Y_2\) are
    both \(E\). Hence the joint \randomvariablesupport is not the product of
    the marginal \randomvariablesupport[supports].
\end{snippetexample}

\end{document}
